\chapter{}
\label{cap:2}

Este cap\'itulo documenta el conjunto de convenciones de ingenier\'ia de
software que rige el desarrollo del proyecto: la organizaci\'on jer\'arquica
de directorios dentro de \texttt{app/}, y las normas de nomenclatura
aplicables a archivos, funciones, variables y constantes en cada uno de los
lenguajes que componen el sistema (TypeScript, Python, MATLAB y C). Estas
normas est\'an codificadas de manera prescriptiva en el archivo
\texttt{app/CONVENTIONS.md}, que act\'ua como fuente de verdad; el presente
cap\'itulo las reproduce, las contextualiza y explica la motivaci\'on de
ingenier\'ia detr\'as de cada una, de modo que puedan aplicarse con criterio
y no como reglas arbitrarias.

\section{Principios generales de dise\~no del c\'odigo}

Con independencia del lenguaje de implementaci\'on, todo el c\'odigo del
repositorio debe adherirse a cuatro principios transversales, orientados a
mantener la base de c\'odigo legible, verificable y f\'acil de modificar a
medida que el proyecto crece.

\begin{description}
    \item[Principio de responsabilidad \'unica por archivo.] Cada archivo de
    c\'odigo fuente debe encapsular una \'unica responsabilidad funcional
    claramente delimitada. Cuando un script o m\'odulo comienza a resolver
    m\'as de una preocupaci\'on independiente (por ejemplo, descargar datos
    \emph{y} transformarlos \emph{y} validarlos), debe descomponerse en
    archivos separados, cada uno responsable de una \'unica etapa. Este
    principio favorece la capacidad de prueba unitaria, reduce el
    acoplamiento entre componentes y facilita la localizaci\'on de defectos.

    \item[L\'imites m\'aximos de extensi\'on por archivo.] Se establecen
    l\'imites expl\'icitos de l\'ineas de c\'odigo por archivo, diferenciados
    seg\'un el lenguaje: $150$ l\'ineas para Python, $100$ para MATLAB, $200$
    para TypeScript y $200$ para C. Estos umbrales no son un fin en s\'i
    mismos, sino un indicador indirecto de violaci\'on del principio de
    responsabilidad \'unica: un archivo que excede el l\'imite es, en la
    pr\'actica totalidad de los casos, un candidato a ser refactorizado en
    unidades m\'as peque\~nas y cohesivas.

    \item[Prohibici\'on de literales num\'ericos sin contexto sem\'antico
    (\emph{magic numbers}).] Todo valor num\'erico que posea un significado
    de negocio o de dominio —umbrales de decisi\'on, tama\~nos m\'aximos de
    bloque, tiempos de espera, n\'umeros de repetici\'on, entre otros— debe
    declararse expl\'icitamente como una constante con nombre descriptivo en
    la cabecera del archivo correspondiente, en lugar de aparecer como un
    literal inline dentro de la l\'ogica de negocio. Esta pr\'actica mejora
    la trazabilidad del significado de cada valor y centraliza su
    modificaci\'on futura en un \'unico punto del c\'odigo.

    \item[Prohibici\'on de c\'odigo fuente comentado (\emph{dead code}).] No
    se permite conservar bloques de c\'odigo inactivo delimitados por
    comentarios como mecanismo de respaldo o referencia hist\'orica. El
    historial de control de versiones (\texttt{git}) cumple esa funci\'on de
    manera nativa y auditable; mantener c\'odigo comentado en el \'arbol de
    trabajo introduce ruido visual, incrementa el costo cognitivo de lectura
    y puede inducir a error respecto de qu\'e ruta de ejecuci\'on es la
    vigente.
\end{description}

\section{Directorio de solo lectura: \texttt{external/}}

El directorio \texttt{app/external/} tiene car\'acter de \textbf{solo
lectura} y no debe, bajo ninguna circunstancia, ser modificado, regenerado o
sobrescrito por c\'odigo del proyecto. Contiene los binarios ejecutables y
las bibliotecas de terceros correspondientes a los solvers de programaci\'on
semidefinida integrados (DSDP, CSDP, SDPA, SDPT3, SeDuMi) y a la herramienta
externa \texttt{InstanceSpace}, distribuidos exactamente en la forma en que
sus respectivos autores los publican.

Esta restricci\'on responde a una consideraci\'on de reproducibilidad e
integridad: los resultados num\'ericos producidos por cada solver dependen
cr\'iticamente de la versi\'on exacta del binario empleado, por lo que
cualquier alteraci\'on no controlada de estos artefactos comprometer\'ia la
validez de las comparaciones de rendimiento documentadas en
\texttt{app/data/benchmark/}. En consecuencia, todo script que necesite
interactuar con estos componentes debe limitarse a \emph{leer} desde
\texttt{external/}, recibiendo la ruta al binario correspondiente como
par\'ametro de entrada (por ejemplo, a trav\'es de \texttt{argparse} en
Python), sin generar artefactos ni archivos temporales dentro de dicho
directorio.

\section{Organizaci\'on jer\'arquica del directorio \texttt{app/}}

La ubicaci\'on de cada archivo dentro del \'arbol de directorios de
\texttt{app/} no responde a un criterio arbitrario, sino a una convenci\'on
expl\'icita que asigna a cada categor\'ia de contenido un directorio
can\'onico \'unico. Esta correspondencia se resume en la
Tabla~\ref{tab:jerarquia-carpetas}.

\begin{table}[htbp]
\centering
\caption{Correspondencia entre tipo de contenido y ubicaci\'on can\'onica
dentro de \texttt{app/}.}
\label{tab:jerarquia-carpetas}
\begin{tabular}{ll}
\toprule
\textbf{Tipo de contenido} & \textbf{Directorio} \\
\midrule
Archivos de configuraci\'on en formato JSON & \texttt{config/} \\
Instancias de programaci\'on semidefinida (\texttt{.dat-s}) & \texttt{data/} \\
Scripts de descarga y preparaci\'on de instancias & \texttt{scripts/instances/} \\
Scripts de extracci\'on de caracter\'isticas (\emph{features}) & \texttt{scripts/features/} \\
Scripts de \emph{benchmarking} de solvers & \texttt{scripts/benchmark/} \\
Scripts de algoritmos gen\'eticos & \texttt{scripts/genetic/} \\
Funciones auxiliares de uso compartido entre m\'odulos & \texttt{scripts/utils/} \\
Scripts exploratorios o de car\'acter desechable & \texttt{scripts/playground/} \\
Binarios de terceros (SDPA, CSDP, DSDP, \ldots) & \texttt{external/} \\
C\'odigo fuente de la aplicaci\'on Electron (main, preload, renderer) & \texttt{src/} \\
\bottomrule
\end{tabular}
\end{table}

Adicionalmente, tanto \texttt{config/} como \texttt{data/} se subdividen en
subdirectorios homónimos que reflejan las etapas conceptuales del pipeline
de an\'alisis introducido en el Cap\'itulo~\ref{cap:1}:
\texttt{instances/}, \texttt{solvers/}, \texttt{benchmark/},
\texttt{features/}, \texttt{metadata/} e \texttt{isa/}. Esta simetr\'ia es
deliberada: la configuraci\'on y los datos generados de una misma etapa del
pipeline residen siempre bajo el mismo nombre de subdirectorio en ambas
ramas del \'arbol, lo que permite localizar de forma inmediata, dado el
nombre de una etapa, tanto sus par\'ametros de entrada como sus artefactos
de salida.

\section{Convenciones de nomenclatura por lenguaje}

Cada lenguaje empleado en el proyecto adopta las convenciones de
nomenclatura est\'andar de su respectivo ecosistema, con el objetivo de que
el c\'odigo resulte idiom\'atico y familiar para desarrolladores con
experiencia previa en dicho lenguaje.

\subsection{Python}
\begin{itemize}
    \item Nombres de archivos y de funciones en \texttt{snake\_case} —por
    ejemplo, \texttt{extract\_features.py} como nombre de archivo y
    \texttt{compute\_sparsity()} como nombre de funci\'on.
    \item Constantes declaradas en \texttt{UPPER\_SNAKE\_CASE} —por ejemplo,
    \texttt{MAX\_BLOCK\_SIZE = 500}.
    \item Clases declaradas en \texttt{PascalCase} —por ejemplo,
    \texttt{InstanceParser}.
\end{itemize}

\subsection{MATLAB}
\begin{itemize}
    \item Nombres de archivos y de funciones en \texttt{camelCase} —por
    ejemplo, \texttt{runSdpt3.m} como nombre de archivo y
    \texttt{parseResult()} como nombre de funci\'on.
    \item Constantes declaradas en \texttt{UPPER\_SNAKE\_CASE}.
\end{itemize}

\subsection{TypeScript}
\begin{itemize}
    \item Nombres de archivos en \texttt{camelCase} —por ejemplo,
    \texttt{renderer.ts}.
    \item Nombres de funciones y variables en \texttt{camelCase}.
    \item Nombres de tipos e interfaces en \texttt{PascalCase}.
\end{itemize}

\subsection{C}
\begin{itemize}
    \item Nombres de archivos en \texttt{snake\_case} —por ejemplo,
    \texttt{sdpa\_wrapper.c}.
    \item Nombres de funciones en \texttt{snake\_case}.
    \item Constantes y macros de preprocesador en \texttt{UPPER\_SNAKE\_CASE}.
\end{itemize}

\section{Convenciones espec\'ificas del lenguaje Python}

Adem\'as de las convenciones de nomenclatura generales, todo script escrito
en Python dentro de este proyecto debe satisfacer los siguientes requisitos
funcionales:

\begin{itemize}
    \item \textbf{Ejecutabilidad aut\'onoma desde l\'inea de comandos.} Cada
    script debe ser invocable de forma independiente y exponer su punto de
    entrada mediante el idioma est\'andar
    \texttt{if \_\_name\_\_ == "\_\_main\_\_":}, evitando dependencias
    impl\'icitas de orquestaci\'on externas al propio script.
    \item \textbf{Parametrizaci\'on mediante \texttt{argparse}.} Ninguna ruta
    de archivo ni par\'ametro de configuraci\'on debe quedar codificado de
    forma fija (\emph{hardcoded}) dentro del cuerpo del script; todos los
    par\'ametros variables se reciben como argumentos de l\'inea de comandos
    procesados mediante el m\'odulo \texttt{argparse} de la biblioteca
    est\'andar.
    \item \textbf{Separaci\'on de canales de salida.} El progreso de
    ejecuci\'on se reporta por \texttt{stdout}, mientras que los mensajes de
    error se reportan exclusivamente por \texttt{stderr}. Esta separaci\'on
    no es meramente estil\'istica: es la que permite, tal como se describe en
    el Cap\'itulo~\ref{cap:1}, que la interfaz de usuario distinga
    program\'aticamente las l\'ineas de progreso normal de las l\'ineas de
    error al renderizar la bit\'acora de ejecuci\'on en tiempo real mediante
    los eventos \texttt{python-stdout} y \texttt{python-stderr}.
    \item \textbf{Convenci\'on de c\'odigos de salida.} El proceso debe
    finalizar con c\'odigo de salida \texttt{0} \'unicamente en caso de
    ejecuci\'on exitosa, y con un c\'odigo distinto de cero ante cualquier
    condici\'on de error, de manera que el proceso principal de Electron
    (evento \texttt{python-exit}) pueda determinar program\'aticamente el
    resultado de la ejecuci\'on sin necesidad de parsear el contenido textual
    de la salida.
\end{itemize}

\section{Convenciones espec\'ificas del lenguaje MATLAB}

\begin{itemize}
    \item Cada archivo con extensi\'on \texttt{.m} debe exponer exactamente
    una funci\'on p\'ublica, cuyo nombre coincide obligatoriamente con el
    nombre del archivo que la contiene, en conformidad con el requisito del
    int\'erprete de MATLAB.
    \item Los resultados producidos por una funci\'on deben retornarse como
    una \'unica estructura (\texttt{struct}) con campos nombrados,
    prefiri\'endose expl\'icitamente esta forma por sobre el retorno de
    m\'ultiples variables de salida posicionales, con el fin de mejorar la
    legibilidad y la estabilidad de la interfaz de la funci\'on ante
    futuras extensiones.
    \item No se invoca la funci\'on \texttt{addpath} dentro de los propios
    scripts de MATLAB: la resoluci\'on de rutas y dependencias se gestiona de
    forma centralizada desde el proceso principal de Electron, evitando que
    la configuraci\'on del \emph{path} de MATLAB quede dispersa entre
    m\'ultiples archivos.
\end{itemize}

\section{Convenciones espec\'ificas del lenguaje TypeScript}

\begin{itemize}
    \item \textbf{Ausencia de estilos en l\'inea.} No se admite el uso de
    atributos de estilo (\texttt{style="..."}) directamente en el marcado
    HTML; toda propiedad visual se expresa mediante clases CSS definidas en
    la hoja de estilos centralizada \texttt{styles.css}, preservando la
    separaci\'on entre estructura y presentaci\'on.
    \item \textbf{Confinamiento de las llamadas de comunicaci\'on entre
    procesos.} Las invocaciones a las APIs de comunicaci\'on entre procesos
    de Electron (\texttt{ipcMain}, \texttt{ipcRenderer}) se restringen
    exclusivamente a los archivos \texttt{main.ts} y \texttt{preload.ts}. El
    c\'odigo del proceso de renderizado no debe, en ning\'un caso, importar
    el m\'odulo \texttt{electron} de forma directa, en conformidad con el
    modelo de seguridad de aislamiento de contexto descrito en el
    Cap\'itulo~\ref{cap:1}.
    \item \textbf{Centralizaci\'on de cadenas de texto traducibles.} Toda
    cadena de texto visible en la interfaz de usuario debe canalizarse a
    trav\'es del mecanismo de internacionalizaci\'on del proyecto, basado en
    el atributo \texttt{data-i18n} y en un diccionario de
    \texttt{translations}; no se admite texto de interfaz codificado de
    forma literal dentro del marcado HTML ni dentro de la l\'ogica
    TypeScript.
\end{itemize}
